% =====================================================
% 题型：圆锥与球表面积之比选择题 (q_solid_cone_sphere_ratio.tex)
% =====================================================
% =====================================================
% 难度：★★ (中档)
% =====================================================
\newcommand{\qSolidConeSphereRatioStem}{据《九章算术》记载，我国匠人常需计算不同几何体表面积或体积的比例以优化用料，例如，制作圆锥形与球形装饰物时，需比较两者的表面积以确定所需涂漆或覆盖材料的用量。若圆锥的底面直径和母线都等于球的直径，则圆锥与球的表面积之比为 \hfill （\quad）}

\newcommand{\qSolidConeSphereRatioOptions}{%
  \mcoptions[4]{\(\dfrac23\)}{\(\dfrac2\pi\)}{\(\dfrac34\)}{\(\dfrac4\pi\)}%
}

\newcommand{\qSolidConeSphereRatioAnswer}{C}

\newcommand{\qSolidConeSphereRatioSolution}{%
  设球的半径为 \(R\)，则球的直径为 \(2R\)。
  
  因为圆锥的底面直径和母线都等于球的直径，所以圆锥的底面半径为 \(r = R\)，圆锥的母线长为 \(l = 2R\)。
  
  圆锥的表面积为底面积与侧面积之和：
  \[
    S_{\text{锥}} = \pi r^2 + \pi rl = \pi R^2 + \pi R \cdot 2R = 3\pi R^2.
  \]
  
  球的表面积公式为：
  \[
    S_{\text{球}} = 4\pi R^2.
  \]
  
  因此，圆锥与球的表面积之比为：
  \[
    \frac{S_{\text{锥}}}{S_{\text{球}}} = \frac{3\pi R^2}{4\pi R^2} = \frac{3}{4}.
  \]
  故选 C。%
}
